\begin{exercisechunk}
\item \textbf{What the universal entropy bound permits.}\suppmark
\label[exercise]{ex:l15-yang-entropy}
\exerciseprofile{3}{2}{3}{\exproof\quad\excalculation\quad\exinterpretation}{%
Use the entropy exponent to distinguish permissible structural adaptation
from an impossible uniform rate on a larger class.}
Suppose a regression procedure has squared-risk success set $\cS_\gamma$ at
rate $n^{-\gamma}$. The entropy theorem gives
\[
 M(\epsilon,\cS_\gamma)
 \lesssim \epsilon^{-2(1-\gamma)/\gamma},
 \qquad 0<\gamma<1.
\]
Assume that the packing entropy of an $r$-dimensional $\alpha$-smooth ball
has order $\epsilon^{-r/\alpha}$.
\begin{enumerate}[(a),beginpenalty=10000]
\item Substitute $\gamma_r=2\alpha/(2\alpha+r)$ into the entropy exponent and
show that the upper bound matches the smoothness-ball entropy order.
\item Consider all functions $x\mapsto g(W^\top x)$, where $g$ ranges over
an $s$-dimensional $\alpha$-smooth ball and $W$ ranges over orthonormal
$d\times s$ matrices. Suppose $d$ is fixed and that covering the possible
orientations contributes $O(ds\log(1/\epsilon))$ to the entropy. Explain why
the entropy theorem permits the rate $n^{-2\alpha/(2\alpha+s)}$ uniformly over
all orientations.
\item Show why the same $s$-dimensional rate cannot hold uniformly over a
full $d$-dimensional smoothness ball when $s<d$.
\item Explain how one procedure can adapt to several unknown smoothness
levels without contradicting the entropy theorem. Your explanation should use
rate-indexed success sets and constants in addition to asymptotic notation.
\end{enumerate}
\end{exercisechunk}
